%\documentstyle[twocolumn,prl,aps]{revtex4}
\documentclass[twocolumn,prl,aps]{revtex4}
%\tightenlines
%\documentstyle[aps,preprint,epsfig]{revtex}
\def\inseps#1#2{\def\epsfsize##1##2{#2##1} \centerline{\epsfbox{#1}}}
\begin{document}
%\baselineskip .6cm

{\bf Comment on ``Dynamic behavior of anisotropic non-equilibrium
driving lattice gases''}


In a recent Letter Albano and Saracco (AS)~\cite{as-02}  study the dynamic
behavior of some anisotropic driven lattice gases by
Monte Carlo (MC) simulations~\cite{STMC}.
They check the scaling of the order parameter $OP$ defined in their Eq.~(4)
and, by data collapse, determine the relevant scaling exponents. Then,
they formulate scaling Ans\"atze (see Eqs.~(5) and (8) 
of AS) in order to relate the measured scaling exponents 
with the critical exponents analytically computed within different 
field-theoretical approaches 
(see Refs.~\cite{js-lc-86,Garrido98}). 
Comparing MC results and theoretical predictions
they conclude that the universality class of the phase 
transition in the driven lattice gases considered in the paper 
is well-described by the field theory proposed in
Ref.~\cite{Garrido98}, ruling out that of Ref.~\cite{js-lc-86}.
In this Comment we point out that the 
Ans\"atze (5) and (8) 
do not take properly into account the strongly anisotropic nature of
the phase transition, by implicitly assuming $z = z_{\bot} = z_{\parallel}$. 
As a consequence, at variance with 
AS's claims, MC data are not conclusive on the issue of the  
universality class.

Let us outline the argument. 
During the relaxation of a spin system from an initial 
uncorrelated state with small magnetization $m_0$, the 
$k$-th moment of the magnetization is expected to scale according 
to~\cite{STMC}
$
M^{(k)}(t,\tau, L, m_0) = 
b^{-k\beta/\nu}f_k(b^{-z}t,b^{1/\nu},b^{1/\nu}\tau,b^{-1}L, b^{x_0}m_0)
$,
where $x_0$ is related to the initial-slip exponent.
The generalization of this expression to a system that is strongly anisotropic 
and has a conserved order parameter requires some care.
Using the notations of Ref.~\cite{sz-95}, one expects
$
M^{(k)}(t,\tau, L_\parallel,L_\bot, m_0) = 
b^{-k\beta/\nu_\parallel}
 f_k(b^{-z_\parallel}t,b^{1/\nu_\parallel}\tau,b^{-1}L_\parallel, 
   b^{-\nu_\bot/\nu_\parallel}L_\bot, b^{x_0}m_0)
$,
where $M$ can be naturally identified with some suitable order parameter,
for instance with the excess density as in AS. 
Since $m_0$ is approximately zero for the initial configurations, 
taking into account that the $OP$ defined by AS is expected
to scale as $\sqrt{M^{(2)}}$~\cite{as-02}, 
one obtains Eq. (5) of AS with the identification $z = z_\parallel$. 
Correctly, in Table I, AS report the predictions $z = 4/3$ and 
$z \approx 1.998$ for the two models they consider (see Eqs. (1) and (2) of AS),
which refer indeed to $z_\parallel$. Formulae 
(6) and (7) follow from Eq. (5) and therefore they are correct only if 
we identify $z = z_\parallel$. 
%%%%%%%%%%%%%%%%%%%%%%%%%%%%%%%%%%%%%%
%%%%%%%%%%%%%%%%%%%%%%%%%%%%%%%%%%%%%%
In Eq.~(8), AS generalize the scaling behavior of $OP$ to
the case of the relaxation 
from the ground state configuration.
Without attemping a derivation of their equation, we simply note
that, as it stands, it is not compatible with Eq.~(5) of AS:
The exponent $z$ appearing in Eq.~(8) must be identified with 
$z_\bot = \nu_\parallel z_\parallel/\nu_\bot$ and not with $z_\parallel$.
Indeed, whatever the initial condition is, we expect that   
after a time interval long enough to allow the relaxation of the initial
condition, $OP$ will display a characteristic long-time behavior,
that should be captured by any reasonable scaling form.
Now, from Eq.~(5) (rewritten as before) 
it is easy to derive that, 
in the thermodynamic limit,
$OP(t,\tau) = t^{-\beta/\nu_\parallel z_\parallel}f_{\rm
FD}(t^{1/\nu_\parallel z_\parallel}\tau)$, while from Eq.~(8)
$OP(t,\tau) = t^{-\beta/\nu_\bot z}f_{\rm
GS}(t^{1/\nu_\bot z}\tau)$.
These two 
equations give the same long-time limit only if in the second case 
we identify $z$ with $z_\bot$, i.e.
$
OP(t,\tau,L_\parallel,L_\bot) = b^{-\beta/\nu_\bot} 
OP(b^{-z_\bot} 
t, b^{1/\nu_\bot}\tau, b^{-\nu_\parallel/\nu_\bot}L_\parallel, 
b^{-1}L_\bot) 
$.
Equations (9) and (10) follow from Eq.~(8) and therefore 
here $z$ must be interpreted as $z_\bot$.
As a consequence, some of the theoretical predictions for the
quantities measured by AS and reported in Tab.~I of Ref.~\cite{as-02},
change according to the following table 
\begin{center}
\begin{tabular}{ccc}
\hline
Model &
\( \beta/\nu_{\bot} z (+) \)&
\( c_{\bot} (+) \)\\
\hline
Eq. (1) in Ref.~\cite{as-02}&
$1/4 $ &
$1/4 $ \\
Eq. (2) in Ref.~\cite{as-02}&
$ \approx 0.13 $&
$ \approx 0.27 $\\
\hline
\end{tabular}
\end{center}
Comparing these corrected predictions with the numerical results 
reported in Tab.~I of AS,
we conclude that the value of $\beta/\nu_{\bot} z$ is in agreement with 
Eq.~(1), which is exactly the opposite conclusion with respect to AS, while
the numerical estimate of $c_{\bot}$ differs from the predictions of 
both field-theoretical models considered in AS's paper.

%%%%%%%%%%%%%%%%
%%%%%%%%%%%%%%%%
\vspace{0.5cm}
\noindent
Sergio Caracciolo,${}^{1}$ Andrea  Gambassi,${}^{2}$ Massimiliano
Gubinelli,${}^{3}$ and Andrea Pelissetto${}^{4}$\\
\noindent
{\small ${}^1$~Dipartimento di Fisica and INFN, Universit\`a di Milano, and
INFM-NEST, I-20133 Milano, Italy\\
${}^2$~Max-Planck-Institut f\"ur Metallforschung
%, Heisenbergstr. 3,
%D-70569, Stuttgart, Germany 
and 
Institut f\"ur Theoretische und Angewandte Physik der Universit\"at
Stuttgart, 
%Pfaffenwaldring 57, 
D-70569 Stuttgart, Germany\\
${}^3$~Dipartimento di Matematica Applicata and INFN, Universit\`a di
Pisa, I-56100 Pisa, Italy\\
${}^4$~Dipartimento di Fisica and INFN, Universit\`a di Roma ``La
Sapienza'', I-00185 Roma, Italy\\ }

%PACS numbers: ?????
\vspace{-1cm}
\begin{thebibliography}{99}

\bibitem{as-02} 
E.~V.~Albano and G.~Saracco,
Phys. Rev. Lett. {\bf 88}, 145701 (2002).
%%CITATION = PRLTA,88,145701;%%

\bibitem{STMC}
Z.~B.~Li, L.~Sch\"ulke, and B.~Zheng,
Phys. Rev. Lett. {\bf 74}, 3396 (1995);
K.~Okano, L.~Sch\"ulke, K.~Yamagishi, and B.~Zheng,
Nucl. Phys. B {\bf 485}, 727 (1997); 
B.~Zheng,
Int. J. Mod. Phys. B {\bf 12}, 1419 (1998);
Physica A {\bf 283}, 80 (2000).
%%CITATION = PRLTA,74,3396;%%
%%CITATION = NUPHA,B485,727;%%
%%CITATION = IMPAE,B12,1419;%%
%%CITATION = PHYSA,A283,80;%%

\bibitem{js-lc-86}
H.~K.~Janssen and B.~Schmittmann,
Z. Phys. B {\bf 64}, 503 (1986).
K.~-t.~Leung and J.~L.~Cardy,
J. Stat. Phys. {\bf 44}, 567 (1986); 
Erratum J. Stat. Phys. {\bf 45}, 1087 (1986).
%%CITATION = ZEPYA,B64,503;%%
%%CITATION = JSTPB,45,1087;%%

\bibitem{Garrido98}
P.~L.~Garrido, F.~de~los~Santos, and M.~A.~Mu{\~n}oz,
Phys. Rev. E {\bf 57}, 752 (1998);
P.~L.~Garrido, M.~A.~Mu{\~n}oz, and  F.~de~los~Santos, and M.~A.~Mu{\~n}oz,
Phys. Rev. E {\bf 61}, R4683 (2000).
%%CITATION = PHRVA,E57,752;%%
%%CITATION = PHRVA,E61,4683;%%

%\bibitem{sz-91}
%B.~Schmittmann and R.~K.~P.~Zia,
%Phys. Rev. Lett. {\bf 66}, 357 (1991).

%\bibitem{sch-93}
%B.~Schmittmann,
%Europhys. Lett. {\bf 24}, 109 (1993).

\bibitem{sz-95}
B.~Schmittmann and R.~K.~P. Zia,
{\sl Statistical mechanics of driven diffusive systems},
in Phase Transitions and Critical Phenomena, edited by C.~Domb
  and J.~L. Lebowitz, vol.~17 (Academic Press, London, 1995).

\end{thebibliography}



\end{document}



